\section{The three do-calculus rules: finite soundness proofs}
\label{sec:annex-p16-do-calculus-soundness}

Use the data and conditional-law notation in
\cref{def:p16-do-calculus-data}. The original rule statement and
regime-based proof idea appear in
\cite[Theorem~3 and Appendix, pp.~677, 686--687]{Pearl1995CausalDiagrams}.
The finite construction and both graph bridges are proved below.

\begin{annexitem}{annex:p16.regime-law}{thm:p16-regime-law}{The finite augmented probability space}
\begin{proof}
The original model has outcome space \(\Omega\), with weights
\(q(\omega)\). Form
\(\widetilde\Omega=\Omega\times\{0,1\}^Z\), assigning
\((\omega,e)\) mass \(q(\omega)2^{-|Z|}\). These weights sum
to one. Each old disturbance is its original function of \(\omega\),
and each new root disturbance is the corresponding bit \(e_i\).
Write \(r_i(u_i)=\mathbb P(U_i=u_i)\). For every old disturbance
tuple \(u\) and every new bit tuple \(e\), summing over compatible
original outcomes gives
\[
 \widetilde{\mathbb P}((U_i)_{i\in N}=u,F=e)
 =\mathbb P((U_i)_{i\in N}=u)2^{-|Z|}
 =\prod_{i\in N}r_i(u_i)\prod_{j\in Z}\tfrac12.
\]
The old and new marginal laws are exactly these factors, so the enlarged
disturbance family is mutually independent. This includes constant old
disturbances and impossible old disturbance tuples.

At the original vertices use the equations of the model after \(X=x\),
except that the equation at \(i\in Z\) has the specified two
branches controlled by \(F_i\). Each branch is total on its finite
input domain. Put the new roots before an original topological order.
Every new edge goes from one of those roots to an original vertex,
so this is a topological order of \(\widetilde K\). The augmented
model therefore has a unique solution at every outcome, by
\cref{thm:p16-causal-solution}.

Fix a bit tuple \(e\). Conditional on \(F=e\), its old-vertex
equations are exactly those of the original model after \(X=x\)
and \(T_e=z_{T_e}\): the indices with bit one use the prescribed
constants and the other equations are retained. Uniqueness identifies
their old-vertex solutions pointwise. Moreover
\(\widetilde{\mathbb P}(F=e)=2^{-|Z|}>0\), and the conditional
mass of each original outcome is
\((q(\omega)2^{-|Z|})/2^{-|Z|}=q(\omega)\).
Summing over the outcomes whose intervention solution lies in \(D\)
proves the claimed joint-event identity. The all-zero and all-one
specializations follow from the definition of \(T_e\). If \(Z\)
is empty, there is one empty bit tuple of mass one and both
specializations are the unchanged baseline model.
\end{proof}
\end{annexitem}

\begin{annexitem}{annex:p16.regime-graph-bridges}{thm:p16-regime-graph-bridges}{Two complete ancestral-moral graph arguments}
\begin{proof}
Use the path/moral equivalence proved in
\cref{thm:p16-moral-separation-bridge}. Empty \(Y\) or \(Z\)
makes both claimed implications immediate from absence of endpoint
paths, so assume both nonempty. Write \(C=X\cup W\).
The new regime vertices are roots with one child each. They add
no directed paths between old vertices.

\paragraph{Exchange.}
Let \(L=K_{\underline Z}\) and
\(H=\operatorname{an}_K(Y\cup Z\cup C)\). Then
\[
 H=\operatorname{an}_L(Y\cup Z\cup C).
\]
For the nontrivial inclusion, take a directed path in \(K\) from
an ancestor to the target set. If it meets \(Z\), stop at its
first such vertex. No edge before that first encounter leaves \(Z\),
so the truncated path survives in \(L\) and still ends in the
target set. If it never meets \(Z\), all its edges survive.
The other inclusion follows because \(L\) is a subgraph of \(K\).

The ancestral vertices for the augmented query with endpoint sets
\(Y,\mathcal F_Z\) and conditioning set \(C\cup Z\) are
exactly \(H\cup\mathcal F_Z\). Original vertices are ancestors
of the old targets exactly as before, and the new roots are themselves
targets. Compare the moral graph on these vertices after deleting
\(C\cup Z\) with the moral graph of \(L\) on \(H\).

Among surviving original vertices, every edge of the first moral
graph is also in the second. Indeed, an original directed edge with
neither endpoint in \(Z\) is unchanged. If two surviving original
vertices are co-parents of a child in \(H\), neither parent is in
\(Z\), so both of those arrows remain in \(L\), preserving
their co-parent edge. Edges involving a deleted \(Z\) vertex
do not need preservation.

The surviving neighbors of a regime vertex \(F_i\) in the augmented
moral graph are precisely the surviving original parents of \(i\):
its edge to its child \(i\in Z\) is deleted, while moralization
joins it to those co-parents. Distinct regime vertices are not adjacent,
since they have distinct children and neither is an original parent.
Any two surviving neighbors of \(F_i\) are joined in the moral
graph of \(L\), because they remain co-parents of \(i\in H\).

Suppose an augmented moral path avoiding \(C\cup Z\) joined
some \(F_i\) to \(Y\). Its first neighbor is an original parent
\(a\) of \(i\), outside \(C\cup Z\). The edge \(a\to i\)
survives in \(L\). Replace the initial vertex \(F_i\) by
\(i\in Z\), retaining this edge to \(a\). Replace any internal
regime vertex by the co-parent edge between its two neighboring old
vertices. All other edges were just shown to be present in the moral
graph of \(L\). The resulting undirected walk joins \(Z\) to
\(Y\) while avoiding \(C\). Erasing repeated portions gives a
path. This contradicts the premise
\(Y\mathrel{\perp_L}Z\mid C\) by the moral-separation bridge.
Thus the augmented moral graph separates the claimed endpoints after
deleting \(C\cup Z\), proving the first implication.

\paragraph{Action deletion and the ancestor exception.}
Put
\[
 A=Z(W)=Z\setminus\operatorname{an}_K(C),\qquad B=Z\setminus A,
 \qquad L=K_{\overline A}.
\]
The equality defining \(A\) follows because \(X\) consists of
roots in \(K\), so no vertex of the disjoint set \(Z\) is an
ancestor of \(X\). Assume \(Y\mathrel{\perp_L}Z\mid C\).
We first prove
\[
 A\cap\operatorname{an}_K(Y\cup C)=\varnothing.
\]
Membership in the ancestors of \(C\) is excluded by the definition
of \(A\). If a member of \(A\) had a directed path to \(Y\)
in \(K\), choose the last vertex of \(A\) on that path.
The portion from it to \(Y\) has no further vertex in \(A\),
so none of its incoming edges is removed in \(L\). It avoids
\(C\), since otherwise its starting vertex would be an ancestor
of \(C\). It is therefore an active path in \(L\), from a
member of \(Z\) to \(Y\), given \(C\). This contradicts
the premise and proves the exclusion.

Let \(H=\operatorname{an}_K(Y\cup C)\). Since \(A\) is
outside \(H\), deleting incoming arrows to \(A\) changes
neither the induced graph on \(H\) nor its ancestry of these
targets. Thus \(H=\operatorname{an}_L(Y\cup C)\).
Every member of \(B\) is an ancestor of \(C\) in both graphs:
its directed path to \(C\) cannot pass through \(A\), whose
vertices are not ancestors of \(C\). Consequently
\(\operatorname{an}_L(Y\cup B\cup C)=H\).
The premise also separates \(B\subseteq Z\) from \(Y\)
given \(C\). Applying the moral bridge to this smaller endpoint
set shows that \(B\) and \(Y\) are separated by \(C\) in
the moral graph of \(K\) on \(H\).

For the augmented query with endpoints \(\mathcal F_Z,Y\) and
conditioning \(C\), the ancestral set is \(H\cup\mathcal F_Z\).
If \(i\in A\), its child is outside \(H\), so \(F_i\)
is isolated in this induced moral graph. If \(i\in B\), the
surviving neighbors of \(F_i\) are \(i\) and the parents of
\(i\) that lie outside \(C\). The vertex \(i\) itself is
outside \(C\), and each such parent is joined to \(i\) by
an original edge. Thus all surviving neighbors are already connected
through \(i\) in the old moral graph with \(C\) deleted.

Any hypothetical path from a regime vertex to \(Y\) can therefore
be converted to an old moral walk from \(B\) to \(Y\). If its
initial \(F_i\) has first neighbor \(u\), start at \(i\) and
use the old edge to \(u\) when \(u\ne i\); when \(u=i\),
start directly there. Replace each internal \(F_j\) by the walk
through its child \(j\) connecting its neighboring old vertices,
omitting a zero-length leg when a neighbor is already \(j\).
All these children lie in
\(B\subseteq H\setminus C\); regime vertices in \(A\) are
isolated and cannot occur on the path. Erasing repetitions would give
an old moral path avoiding \(C\), contradicting the separation
of \(B\) and \(Y\). The augmented separation follows. This
proves the second implication with exactly the stated exception.
\end{proof}
\end{annexitem}

\begin{annexitem}{annex:p16.do-rule-one}{thm:p16-do-rule-one}{Observation deletion in the baseline intervention model}
\begin{proof}
Equation replacement at \(X\) leaves an acyclic model on
\(G_{\overline X}\) with the same independent disturbances.
By \cref{thm:p16-finite-global-markov}, the graph condition gives
conditional independence of \(V_Y^{X=x}\) and \(V_Z^{X=x}\)
given \(V_X^{X=x},V_W^{X=x}\). The first conditioning tuple
is constantly \(x\). Since the event \((z,w)\) has positive
mass, its containing event \(w\) does too. Conditional independence
on that stratum gives
\[
 p_x(y,z\mid w)=p_x(y\mid w)p_x(z\mid w).
\]
The factor \(p_x(z\mid w)=p_x(z,w)/p_x(w)\) is positive.
Dividing by it proves the rule. If \(Y\) or \(Z\) is empty,
the same calculation uses its unique tuple and probability one.
\end{proof}
\end{annexitem}

\begin{annexitem}{annex:p16.do-rule-two}{thm:p16-do-rule-two}{Exchange by comparing two positive regime strata}
\begin{proof}
In the baseline model after \(X=x\), consider the event
\(V_Z^{X=x}=z,V_W^{X=x}=w\), whose mass is \(p_x(z,w)>0\).
On this event, additionally replacing the equations at \(Z\) by
their already attained values \(z\) leaves the entire solution
unchanged, by pointwise consistency in
\cref{thm:p16-causal-intervention-laws}. In particular its outcomes
also satisfy \(V_W^{X=x,Z=z}=w\). Hence
\[
 p_{x,z}(w)\geq p_x(z,w)>0.
\]
The sequential and simultaneous replacements agree because \(X,Z\)
are disjoint. This step applies to the full set \(Z\) simultaneously.

Construct the regime model at these fixed values \(x,z\). The
first implication of \cref{thm:p16-regime-graph-bridges} and the
finite global Markov theorem give
\[
 \widetilde V_Y\mathrel{\perp\!\!\!\perp}F
       \mid(\widetilde V_X,\widetilde V_W,\widetilde V_Z).
\]
The joint events consisting of \((x,w,z)\) and the all-zero or
all-one regime have respectively probabilities
\(2^{-|Z|}p_x(z,w)\) and \(2^{-|Z|}p_{x,z}(w)\).
Both are positive; in the all-one regime the coordinates at \(Z\)
are constantly \(z\).

For any finite conditional-independence identity, conditioning further
on a regime value whose joint conditioning event is positive preserves
the response conditional: divide the product identity by the positive
conditional mass of that regime. Thus the two displayed regime strata
give the same conditional distribution of \(\widetilde V_Y\).
The regime-law lemma identifies them respectively as
\(p_x(y\mid z,w)\) and \(p_{x,z}(y\mid w)\), proving the
claim. This proof requires neither independence of variables inside
\(Z\) nor that its vertices be incomparable in the DAG.
\end{proof}
\end{annexitem}

\begin{annexitem}{annex:p16.do-rule-three}{thm:p16-do-rule-three}{Action deletion on the common positive conditioning strata}
\begin{proof}
Use the same regime construction. The second implication of
\cref{thm:p16-regime-graph-bridges} and the finite global Markov
theorem give
\[
 \widetilde V_Y\mathrel{\perp\!\!\!\perp}F
             \mid(\widetilde V_X,\widetilde V_W).
\]
The conditioning event \((x,w)\) together with the all-zero
regime has mass \(2^{-|Z|}p_x(w)>0\), and with the all-one
regime it has mass \(2^{-|Z|}p_{x,z}(w)>0\). As in the preceding
proof, divide conditional independence by the positive conditional
mass of each regime. The response conditional is the same in both.
By \cref{thm:p16-regime-law}, the two conditionals are exactly
\(p_x(y\mid w)\) and \(p_{x,z}(y\mid w)\).

The graph-bridge proof used the set
\(Z\setminus\operatorname{an}_{G_{\overline X}}(W)\).
It did not delete incoming arrows to every member of \(Z\).
The resulting rule consequently permits interventions that affect
\(W\), provided the precise separation holds and both probabilities
of the conditioning event are positive. When \(W\) is empty,
both such probabilities are one. No source of dependence among the
original disturbances was introduced in any of these arguments.
\end{proof}
\end{annexitem}

\begin{annexitem}{annex:p16.do-rule-sequence}{thm:p16-do-rule-sequence}{Every graph license and numerical factor in the chain}
\begin{proof}
The four independent disturbance masses define a finite model with
topological order \(A,B,C,D\). Its full law is
\[
 p(a,b,c,d)=\tfrac14\alpha(b\oplus a)\beta(c\oplus b),
\]
because the unique disturbance tuple giving a specified assignment is
\((a,b\oplus a,c\oplus b,d)\). All these masses are positive.
After fixing \(A=a\), its law on the other coordinates is
\[
 p_a(b,c,d)=\tfrac12\alpha(b\oplus a)\beta(c\oplus b).
\]
In particular \(p_a(b,d)=\frac12\alpha(b\oplus a)>0\),
and \(p_a(b)=\alpha(b\oplus a)>0\).

For the first step, Rule 1 uses baseline intervention \(A=a\),
response \(C\), observation to delete \(D\), and retained
observation \(B=b\). The isolated vertex \(D\) is separated
from \(C\) in the required graph, with conditioning \(\{A,B\}\).
The just-computed probability supplies its support condition.

For the second step, Rule 2 exchanges observation \(B=b\) with
intervention at \(B\), retaining baseline \(A=a\) and no
other observation. Deleting outgoing arrows from \(B\) removes
\(B\to C\), leaving \(C\) isolated from \(B\). This proves
the required separation given \(A\). Its support is \(p_a(b)>0\).

For the third step, take the baseline intervention to be \(B=b\)
and the action to delete to be \(A=a\). There is no conditioning
set \(W\), so the exception set is the entire singleton \(\{A\}\).
Deleting incoming arrows at \(B\) removes \(A\to B\), and
\(A\) is then isolated from \(C\), proving the Rule 3 premise
given \(B\). Both empty conditioning events have probability one.
The order of the two preceding interventions is immaterial by
\cref{thm:p16-causal-intervention-laws}.

The fourth step uses Rule 2 with empty baseline, action \(B=b\),
response \(C\), and no other observations. Deleting outgoing
arrows from \(B\) again isolates \(C\), so the graph test holds.
Summing the original joint law gives
\[
 p(B=b)=\tfrac12\sum_{a=0}^1\alpha(b\oplus a)=\tfrac12>0.
\]
This supplies the last support condition. Finally summing the joint
law over \(A,D\) gives
\(p(B=b,C=c)=\frac12\beta(c\oplus b)\), hence
\(p(c\mid b)=\beta(c\oplus b)\). These calculations verify
all steps and the claimed response-one values directly.
\end{proof}
\end{annexitem}

\begin{annexitem}{annex:p16.do-failed-conditions}{thm:p16-do-failed-conditions}{A wrong ancestor deletion and an undefined conditional}
\begin{proof}
For the first model the outcome weights
\(\frac12\alpha(a)\alpha(b)\) sum to one and factor, so the
coordinate disturbances and the additional constant disturbance are
mutually independent. The equations have topological order
\(U,Z,W,Y\). Their eight full assignments have the following
masses, obtained from \(a=z\oplus u\), \(b=w\oplus z\),
and \(y=u\):
\[
\begin{array}{cccc|c}
 U&Z&W&Y&\text{mass}\\\hline
 0&0&0&0&9/32\\
 0&0&1&0&3/32\\
 0&1&0&0&1/32\\
 0&1&1&0&3/32\\
 1&0&0&1&3/32\\
 1&0&1&1&1/32\\
 1&1&0&1&3/32\\
 1&1&1&1&9/32
\end{array}
\]
Every other assignment has mass zero. Summing the four rows with
\(W=1\) gives \(16/32=1/2\); summing those also having
\(Y=1\) gives \(10/32=5/16\). Thus
\(p(Y=1\mid W=1)=(5/16)/(1/2)=5/8\).

After fixing \(Z=z\), the unchanged root \(U\) remains fair,
and \(W=z\oplus U_W\) depends only on its independent disturbance.
Therefore
\[
 p_z(W=1)=\alpha(1\oplus z)>0,\qquad
 p_z(Y=1,W=1)=\tfrac12\alpha(1\oplus z),
\]
so the corresponding response conditional is \(1/2\). Here the
baseline intervention set is empty. The action vertex \(Z\) is
an ancestor of the conditioned \(W\), making \(Z(W)\) empty.
The correct Rule 3 graph is consequently the original graph, where
\(Z\leftarrow U\to Y\) is active given \(W\). The rule
does not authorize action deletion.

If all incoming arrows to \(Z\) were incorrectly deleted, the
edge \(U\to Z\) would disappear. The remaining graph would
have one component \(Z\to W\) and another \(U\to Y\).
It separates \(Z\) from \(Y\) even after conditioning on \(W\),
so that incorrect test would assert the disproved equality. This
establishes why the exact ancestor exception matters even when both
conditioning events are positive.

In the second model, the sole outcome has mass one and both disturbances
are constants. They are mutually independent, and the graph \(Z\to Y\)
is acyclic. Observationally \((Z,Y)=(0,0)\). Under the intervention
\(Z=1\), the retained response equation gives \(Y=1\), proving
\(p_1(Y=1)=1\). The graph after deleting the outgoing \(Z\)
edge consists of two isolated vertices, so the graphical condition
of Rule 2 holds. Its required probability \(p(Z=1)\), however,
is zero. The observational conditional ratio would have numerator
zero and denominator zero and is not defined by the finite probability
convention. The rule therefore makes no assertion at that value.

For the last model, take two independent fair disturbance bits for
\(Z,Y\) and a constant disturbance for \(W\). The four outcomes
give \((Z,Y)\) mass \(1/4\) at each pair and set \(W=Z\).
Thus \(p(W=1)=1/2\), and the joint event \(Y=1,W=1\) has
mass \(1/4\), giving conditional probability \(1/2\). The
response vertex \(Y\) is isolated, so it is separated from \(Z\)
given \(W\) in the correct Rule 3 graph. Here \(Z(W)\) is
empty because \(Z\to W\). Fixing \(Z=0\) forces \(W=0\)
at every outcome, so the intervention-side event \(W=1\) has
mass zero and its response conditional is undefined. The second
support assumption in Rule 3 is therefore not supplied by the first
one and its graph test.
\end{proof}
\end{annexitem}
